% Zdroj: PDF priklady-jaro-2024.pdf, fyzická str. 21-22. Značka: specializace UI-SU. Zařazeno: S3.
\section{Lineární regrese a regularizace}
\szzotazka{jaro 2024}{27}{Lineární regrese a regularizace}

\noindent\textbf{Zadání.}\par
Chceme natrénovat lineární model tvaru $y = \beta_1 x_1 + \beta_2 x_2 + \beta_0$ na datech níže.
Při trénování chceme použít gradientní metodu a L2 regularizaci. Použijeme učící konstantu
$\alpha = 0.1$ a regularizační konstantu $\lambda = 0.5$.

\begin{center}
\begin{tabular}{cc|c}
\toprule
$x_1$ & $x_2$ & $y$ \\
\midrule
0 & 2 & 3 \\
2 & 3 & 9 \\
1 & 1 & 2 \\
2 & 0 & 2 \\
\bottomrule
\end{tabular}
\end{center}

\begin{enumerate}
  \item Definujte chybovou funkci pro model zadaný podle parametrů nahoře (součet čtverců
    s L2 regularizačním členem).
  \item Předpokládejme, že aktuální odhad koeficientů modelu je $\beta_1 = 1$, $\beta_2 = 3$,
    $\beta_0 = -1$. Proveďte jeden krok gradientní metody pro nastavení koeficientů. Jaké jsou
    jejich nové hodnoty?
\end{enumerate}

\medskip
\noindent\textbf{Náčrt řešení.}\par
\begin{enumerate}
  \item Chybová funkce je typicky $\mathrm{RSS} = \sum_{i=1}^{N} (y_i - \hat{y}_i)^2$, pro L2
    regularizaci přičteme ještě člen $\lambda \sum_{i=1}^{m} \beta_i^2$ ($N$ je počet trénovacích
    vzorů, $m+1$ je počet koeficientů modelu, $\hat{y}_i$ je výstup modelu pro $i$-tý vstup).
    Koeficient $\beta_0$ se většinou v regularizaci neuvažuje.
  \item Je potřeba spočítat gradient chybové funkce podle jednotlivých parametrů. Spočítáme ho
    pro jeden trénovací vzor, pro všechny vzory to bude jen součet. Derivace
    $\frac{d\,\mathrm{RSS}}{d\beta_0} = -2(y_i - \hat{y}_i)$, derivace
    $\frac{d\,\mathrm{RSS}}{d\beta_i} = -2(y_i - \hat{y}_i) x_i$. Nesmíme zapomenout také na
    derivaci regularizačního členu podle $\beta_i$ -- ta je $\lambda 2 \beta_i = \beta_i$.
    Dosazením dostaneme, že derivace RSS podle $\beta_0$ je $4 + 2 + 2 + (-2) = 6$, derivace
    podle $\beta_1$ je 2 a derivace podle $\beta_2$ je 16. Přičtením regularizačního členu
    dostaneme, že derivace chybové funkce podle $\beta_1$ je 3 a podle $\beta_2$ je 19. Hodnoty
    v dalším kroku tedy jsou $\beta_0 = \beta_0 - \alpha \cdot 6 = -1.6$,
    $\beta_1 = \beta_1 - \alpha \cdot 3 = 0.7$ a $\beta_2 = \beta_2 - \alpha \cdot 19 = 1.1$.
\end{enumerate}
